\title{2020 年高考数学 全国III卷-文科}
\date{}
\maketitle

\section*{一、选择题}
本题共12小题，每小题5分，共60分。在每小题给出的四个选项中，只有一项是符合题目要求的。

\begin{question}[points=5]
已知集合 $A=\{1,2,3,5,7,11\}$，$B=\{x|3<x<15\}$，则 $A\cap B$ 中元素的个数为
\begin{choices}
  \item $2$
  \item $3$
  \item $4$
  \item $5$
\end{choices}
\begin{solution}
\textbf{答案：}B

\textbf{分析：}采用列举法列举出 $A\cap B$ 中元素的即可.

\textbf{详解：}由题意，$A\cap B=\{5,7,11\}$，故 $A\cap B$ 中元素的个数为3. 故选：B.
\end{solution}
\end{question}

\begin{question}[points=5]
若 $z(1+i)=1-i$，则 $z=$
\begin{choices}
  \item $1-i$
  \item $1+i$
  \item $-i$
  \item $i$
\end{choices}
\begin{solution}
\textbf{答案：}D

\textbf{分析：}先利用除法运算求得 $z$，再利用共轭复数的概念得到 $z$ 即可.

\textbf{详解：}因为 $z=\frac{1-i}{1+i}=\frac{(1-i)^2}{(1+i)(1-i)}=\frac{-2i}{2}=-i$，所以 $\bar{z}=i$. 故选：D.
\end{solution}
\end{question}

\begin{question}[points=5]
设一组样本数据 $x_1,x_2,\dots,x_n$ 的方差为 $0.01$，则数据 $10x_1,10x_2,\dots,10x_n$ 的方差为
\begin{choices}
  \item $0.01$
  \item $0.1$
  \item $1$
  \item $10$
\end{choices}
\begin{solution}
\textbf{答案：}C

\textbf{分析：}根据新数据与原数据关系确定方差关系，即得结果.

\textbf{详解：}因为数据 $ax_i+b\ (i=1,2,\dots,n)$ 的方差是数据 $x_i\ (i=1,2,\dots,n)$ 的方差的 $a^2$ 倍，所以所求数据方差为 $10^2\times0.01=1$. 故选：C.
\end{solution}
\end{question}

\begin{question}[points=5]
Logistic 模型是常用数学模型之一，可应用于流行病学领域. 有学者根据公布数据建立了某地区新冠肺炎累计确诊病例数 $I(t)$ ($t$ 的单位：天)的 Logistic 模型：$I(t)=\frac{K}{1+e^{-0.23(t-53)}}$，其中 $K$ 为最大确诊病例数. 当 $I(t^*)=0.95K$ 时，标志着已初步遏制疫情，则 $t^*$ 约为 ( $\ln19\approx3$ )
\begin{choices}
  \item $60$
  \item $63$
  \item $66$
  \item $69$
\end{choices}
\begin{solution}
\textbf{答案：}C

\textbf{分析：}将 $t=t^*$ 代入函数 $I(t)=\frac{K}{1+e^{-0.23(t-53)}}$ 结合 $I(t^*)=0.95K$ 求得 $t^*$ 即可得解.

\textbf{详解：}因为 $I(t)=\frac{K}{1+e^{-0.23(t-53)}}$，所以 $I(t^*)=\frac{K}{1+e^{-0.23(t^*-53)}}=0.95K$，则 $e^{0.23(t^*-53)}=19$，所以 $0.23(t^*-53)=\ln19\approx3$，解得 $t^*\approx\frac{3}{0.23}+53\approx66$. 故选：C.
\end{solution}
\end{question}

\begin{question}[points=5]
已知 $\sin\theta+\sin(\theta+\frac{\pi}{3})=1$，则 $\sin(\theta+\frac{\pi}{6})=$
\begin{choices}
  \item $\frac{1}{2}$
  \item $\frac{\sqrt{3}}{3}$
  \item $\frac{\sqrt{2}}{3}$
  \item $\frac{\sqrt{2}}{2}$
\end{choices}
\begin{solution}
\textbf{答案：}B

\textbf{分析：}将所给的三角函数式展开变形，然后再逆用两角和的正弦公式即可求得三角函数式的值.

\textbf{详解：}由题意可得：$\sin\theta+\frac{1}{2}\sin\theta+\frac{\sqrt{3}}{2}\cos\theta=1$，则 $\frac{3}{2}\sin\theta+\frac{\sqrt{3}}{2}\cos\theta=1$，即 $\frac{\sqrt{3}}{2}\sin\theta+\frac{1}{2}\cos\theta=\frac{\sqrt{3}}{3}$，从而有 $\sin\theta\cos\frac{\pi}{6}+\cos\theta\sin\frac{\pi}{6}=\frac{\sqrt{3}}{3}$，即 $\sin(\theta+\frac{\pi}{6})=\frac{\sqrt{3}}{3}$. 故选：B.
\end{solution}
\end{question}

\begin{question}[points=5]
在平面内，$A$，$B$ 是两个定点，$C$ 是动点，若 $\overrightarrow{AC}\cdot\overrightarrow{BC}=1$，则点 $C$ 的轨迹为
\begin{choices}
  \item 圆
  \item 椭圆
  \item 抛物线
  \item 直线
\end{choices}
\begin{solution}
\textbf{答案：}A

\textbf{分析：}首先建立平面直角坐标系，然后结合数量积的定义求解其轨迹方程即可.

\textbf{详解：}设 $|AB|=2a\ (a>0)$，以 $AB$ 中点为坐标原点建立平面直角坐标系，则 $A(-a,0),B(a,0)$，设 $C(x,y)$，可得 $\overrightarrow{AC}=(x+a,y),\overrightarrow{BC}=(x-a,y)$，从而 $\overrightarrow{AC}\cdot\overrightarrow{BC}=(x+a)(x-a)+y^2$，结合题意得 $(x+a)(x-a)+y^2=1$，整理得 $x^2+y^2=a^2+1$，即点 $C$ 的轨迹是以 $AB$ 中点为圆心，$\sqrt{a^2+1}$ 为半径的圆. 故选：A.
\end{solution}
\end{question}

\begin{question}[points=5]
设 $O$ 为坐标原点，直线 $x=2$ 与抛物线 $C:y^2=2px(p>0)$ 交于 $D$，$E$ 两点，若 $OD\perp OE$，则 $C$ 的焦点坐标为
\begin{choices}
  \item $(\frac{1}{4},0)$
  \item $(\frac{1}{2},0)$
  \item $(1,0)$
  \item $(2,0)$
\end{choices}
\begin{solution}
\textbf{答案：}B

\textbf{分析：}根据题中所给的条件 $OD\perp OE$，结合抛物线的对称性，可知 $\angle DOx=\angle COx=\frac{\pi}{4}$，从而可以确定出点 $D$ 的坐标，代入方程求得 $p$ 的值，进而求得其焦点坐标.

\textbf{详解：}因为直线 $x=2$ 与抛物线 $y^2=2px(p>0)$ 交于 $C,D$ 两点，且 $OD\perp OE$，根据抛物线的对称性可以确定 $\angle DOx=\angle COx=\frac{\pi}{4}$，所以 $C(2,2)$，代入抛物线方程得 $4=4p$，解得 $p=1$，所以其焦点坐标为 $(\frac{1}{2},0)$. 故选：B.
\end{solution}
\end{question}

\begin{question}[points=5]
点 $(0,-1)$ 到直线 $y=k(x+1)$ 距离的最大值为
\begin{choices}
  \item $1$
  \item $\sqrt{2}$
  \item $\sqrt{3}$
  \item $2$
\end{choices}
\begin{solution}
\textbf{答案：}B

\textbf{分析：}首先根据直线方程判断出直线过定点 $P(-1,0)$，设 $A(0,-1)$，当直线 $y=k(x+1)$ 与 $AP$ 垂直时，点 $A$ 到直线的距离最大，即可求得结果.

\textbf{详解：}由 $y=k(x+1)$ 可知直线过定点 $P(-1,0)$，设 $A(0,-1)$，当直线 $y=k(x+1)$ 与 $AP$ 垂直时，点 $A$ 到直线的距离最大，即为 $|AP|=\sqrt{2}$. 故选：B.
\end{solution}
\end{question}

\begin{question}[points=5]
下图为某几何体的三视图，则该几何体的表面积是
\begin{choices}
  \item $6+4\sqrt{2}$
  \item $4+4\sqrt{2}$
  \item $6+2\sqrt{3}$
  \item $4+2\sqrt{3}$
\end{choices}
\begin{solution}
\textbf{答案：}C

\textbf{分析：}根据三视图特征，在正方体中截取出符合题意的立体图形，求出每个面的面积，即可求得其表面积.

\textbf{详解：}根据三视图特征，在正方体中截取出符合题意的立体图形. 根据立体图形可得 $S_{\triangle ABC}=S_{\triangle ADC}=S_{\triangle CDB}=\frac{1}{2}\times2\times2=2$，根据勾股定理可得 $AB=AD=DB=2\sqrt{2}$，所以 $\triangle ADB$ 是边长为 $2\sqrt{2}$ 的等边三角形，其面积为 $\frac{1}{2}\times(2\sqrt{2})^2\times\frac{\sqrt{3}}{2}=2\sqrt{3}$. 故该几何体的表面积为 $3\times2+2\sqrt{3}=6+2\sqrt{3}$. 故选：C.
\end{solution}
\end{question}

\begin{question}[points=5]
设 $a=\log_32$，$b=\log_53$，$c=\frac{2}{3}$，则
\begin{choices}
  \item $a<c<b$
  \item $a<b<c$
  \item $b<c<a$
  \item $c<a<b$
\end{choices}
\begin{solution}
\textbf{答案：}A

\textbf{分析：}分别将 $a,b$ 改写为 $a=\frac{1}{3}\log_32^3$，$b=\frac{1}{3}\log_53^3$，再利用单调性比较即可.

\textbf{详解：}因为 $a=\frac{1}{3}\log_32^3<\frac{1}{3}\log_39=\frac{2}{3}=c$，$b=\frac{1}{3}\log_53^3>\frac{1}{3}\log_525=\frac{2}{3}=c$，所以 $a<c<b$. 故选：A.
\end{solution}
\end{question}

\begin{question}[points=5]
在 $\triangle ABC$ 中，$\cos C=\frac{2}{3}$，$AC=4$，$BC=3$，则 $\tan B=$
\begin{choices}
  \item $\sqrt{5}$
  \item $2\sqrt{5}$
  \item $4\sqrt{5}$
  \item $8\sqrt{5}$
\end{choices}
\begin{solution}
\textbf{答案：}C

\textbf{分析：}先根据余弦定理求 $c$，再根据余弦定理求 $\cos B$，最后根据同角三角函数关系求 $\tan B$.

\textbf{详解：}设 $AB=c, BC=a=3, CA=b=4$，由余弦定理得 $c^2=a^2+b^2-2ab\cos C=9+16-2\times3\times4\times\frac{2}{3}=9$，所以 $c=3$，则 $\cos B=\frac{a^2+c^2-b^2}{2ac}=\frac{9+9-16}{2\times3\times3}=\frac{1}{9}$，$\sin B=\sqrt{1-(\frac{1}{9})^2}=\frac{4\sqrt{5}}{9}$，所以 $\tan B=4\sqrt{5}$. 故选：C.
\end{solution}
\end{question}

\begin{question}[points=5]
已知函数 $f(x)=\sin x+\frac{1}{\sin x}$，则
\begin{choices}
  \item $f(x)$ 的最小值为 $2$
  \item $f(x)$ 的图像关于 $y$ 轴对称
  \item $f(x)$ 的图像关于直线 $x=\pi$ 对称
  \item $f(x)$ 的图像关于直线 $x=\frac{\pi}{2}$ 对称
\end{choices}
\begin{solution}
\textbf{答案：}D

\textbf{分析：}根据基本不等式使用条件可判断A；根据奇偶性可判断B；根据对称性判断C,D.

\textbf{详解：}因为 $\sin x$ 可以为负，所以A错；$\sin x\neq0$，$x\neq k\pi(k\in\mathbb Z)$，$f(-x)=-\sin x-\frac{1}{\sin x}=-f(x)$，所以 $f(x)$ 关于原点对称；$f(2\pi-x)=-\sin x-\frac{1}{\sin x}\neq f(x)$，$f(\pi-x)=\sin x+\frac{1}{\sin x}=f(x)$，所以 $f(x)$ 关于直线 $x=\frac{\pi}{2}$ 对称，故C错，D对. 故选：D.
\end{solution}
\end{question}

\section*{二、填空题}
本题共4小题，每小题5分，共20分。

\begin{question}[points=5]
若 $x$，$y$ 满足约束条件 $\begin{cases} x+y\geq0 \\ 2x-y\geq0 \\ x\leq1 \end{cases}$，则 $z=3x+2y$ 的最大值为
\fillin[{$7$}]
\begin{solution}
\textbf{答案：}$7$

\textbf{分析：}作出可行域，利用截距的几何意义解决.

\textbf{详解：}不等式组所表示的可行域如图. 因为 $z=3x+2y$，所以 $y=-\frac{3}{2}x+\frac{z}{2}$，易知截距 $\frac{z}{2}$ 越大，则 $z$ 越大. 平移直线 $y=-\frac{3}{2}x$，当 $y=-\frac{3}{2}x+\frac{z}{2}$ 经过 $A$ 点时截距最大，此时 $z$ 最大. 由 $\begin{cases} y=2x \\ x=1 \end{cases}$ 得 $\begin{cases} x=1 \\ y=2 \end{cases}$，$A(1,2)$，所以 $z_{\max}=3\times1+2\times2=7$.
\end{solution}
\end{question}

\begin{question}[points=5]
设双曲线 $C:\frac{x^2}{a^2}-\frac{y^2}{b^2}=1\ (a>0,b>0)$ 的一条渐近线为 $y=\sqrt{2}x$，则 $C$ 的离心率为
\fillin[{$\sqrt{3}$}]
\begin{solution}
\textbf{答案：}$\sqrt{3}$

\textbf{分析：}根据已知可得 $\frac{b}{a}=\sqrt{2}$，结合双曲线中 $a,b,c$ 的关系，即可求解.

\textbf{详解：}由双曲线方程可得其焦点在 $x$ 轴上，因为其一条渐近线为 $y=\sqrt{2}x$，所以 $\frac{b}{a}=\sqrt{2}$，$e=\frac{c}{a}=\sqrt{1+\frac{b^2}{a^2}}=\sqrt{3}$.
\end{solution}
\end{question}

\begin{question}[points=5]
设函数 $f(x)=\frac{e^x}{x+a}$，若 $f'(1)=\frac{e}{4}$，则 $a=$
\fillin[{$1$}]
\begin{solution}
\textbf{答案：}$1$

\textbf{分析：}由题意首先求得导函数的解析式，然后得到关于实数 $a$ 的方程，解方程即可确定实数 $a$ 的值.

\textbf{详解：}由函数的解析式可得 $f'(x)=\frac{e^x(x+a)-e^x}{(x+a)^2}=\frac{e^x(x+a-1)}{(x+a)^2}$，则 $f'(1)=\frac{e^1\times(1+a-1)}{(1+a)^2}=\frac{ae}{(a+1)^2}$，据此可得 $\frac{ae}{(a+1)^2}=\frac{e}{4}$，整理得 $a^2-2a+1=0$，解得 $a=1$.
\end{solution}
\end{question}

\begin{question}[points=5]
已知圆锥的底面半径为 $1$，母线长为 $3$，则该圆锥内半径最大的球的体积为
\fillin[{$\frac{\sqrt{2}}{3}\pi$}]
\begin{solution}
\textbf{答案：}$\frac{\sqrt{2}}{3}\pi$

\textbf{分析：}将原问题转化为求解圆锥内切球的问题，然后结合截面确定其半径即可确定体积的值.

\textbf{详解：}易知半径最大球为圆锥的内切球，球与圆锥内切时的轴截面如图所示，其中 $BC=2,AB=AC=3$，且点 $M$ 为 $BC$ 边上的中点，设内切圆的圆心为 $O$. 由于 $AM=\sqrt{3^2-1^2}=2\sqrt{2}$，故 $S_{\triangle ABC}=\frac{1}{2}\times2\times2\sqrt{2}=2\sqrt{2}$. 设内切圆半径为 $r$，则 $S_{\triangle ABC}=S_{\triangle AOB}+S_{\triangle BOC}+S_{\triangle AOC}=\frac{1}{2}\times AB\times r+\frac{1}{2}\times BC\times r+\frac{1}{2}\times AC\times r=\frac{1}{2}\times(3+3+2)\times r=2\sqrt{2}$，解得 $r=\frac{\sqrt{2}}{2}$，其体积 $V=\frac{4}{3}\pi r^3=\frac{\sqrt{2}}{3}\pi$.
\end{solution}
\end{question}

\section*{三、解答题}
共70分。解答应写出文字说明、证明过程或演算步骤。第17～21题为必考题，每个试题考生都必须作答。第22、23题为选考题，考生根据要求作答。

\begin{question}[points=12]
设等比数列 $\{a_n\}$ 满足 $a_1+a_2=4$，$a_3-a_1=8$.

(1) 求 $\{a_n\}$ 的通项公式；

(2) 记 $S_n$ 为数列 $\{\log_3 a_n\}$ 的前 $n$ 项和，若 $S_m+S_{m+1}=S_{m+3}$，求 $m$.
\begin{solution}
\textbf{答案：}(1) $a_n=3^{n-1}$ (2) $m=6$

\textbf{分析：}(1)设等比数列 $\{a_n\}$ 的公比为 $q$，根据题意列出方程组，求得首项和公比，进而求得通项公式；(2)由(1)求出 $\{\log_3 a_n\}$ 的通项公式，利用等差数列求和公式求得 $S_n$，根据已知列出关于 $m$ 的等量关系式，求得结果.

\textbf{详解：}(1) 设等比数列 $\{a_n\}$ 的公比为 $q$，根据题意有 $\begin{cases} a_1+a_1q=4 \\ a_1q^2-a_1=8 \end{cases}$，解得 $\begin{cases} a_1=1 \\ q=3 \end{cases}$，所以 $a_n=3^{n-1}$.
(2) 令 $b_n=\log_3 a_n=\log_3 3^{n-1}=n-1$，则 $S_n=\frac{n(0+n-1)}{2}=\frac{n(n-1)}{2}$. 由 $S_m+S_{m+1}=S_{m+3}$ 得 $\frac{m(m-1)}{2}+\frac{m(m+1)}{2}=\frac{(m+2)(m+3)}{2}$，整理得 $m^2-5m-6=0$，因为 $m>0$，所以 $m=6$.
\end{solution}
\end{question}

\begin{question}[points=12]
某学生兴趣小组随机调查了某市100天中每天的空气质量等级和当天到某公园锻炼的人次，整理数据得到下表(单位：天)：



| 锻炼人次 | [0,200] | (200,400] | (400,600] |

|----------|---------|-----------|-----------|

| 1(优)    | 2       | 16        | 25        |

| 2(良)    | 5       | 10        | 12        |

| 3(轻度污染) | 6    | 7         | 8         |

| 4(中度污染) | 7    | 2         | 0         |



(1) 分别估计该市一天的空气质量等级为1,2,3,4的概率；

(2) 求一天中到该公园锻炼的平均人次的估计值(同一组中的数据用该组区间的中点值为代表)；

(3) 若某天的空气质量等级为1或2，则称这天“空气质量好”；若某天的空气质量等级为3或4，则称这天“空气质量不好”。根据所给数据，完成下面的 $2\times2$ 列联表，并根据列联表，判断是否有95%的把握认为一天中到该公园锻炼的人次与该市当天的空气质量有关？



附：$K^2=\frac{n(ad-bc)^2}{(a+b)(c+d)(a+c)(b+d)}$，$P(K^2\geq3.841)\approx0.050$.
\begin{solution}
\textbf{答案：}(1) 概率分别为 $0.43$, $0.27$, $0.21$, $0.09$ (2) $350$ (3) 有

\textbf{分析：}(1)根据频数分布表可计算出该市一天的空气质量等级分别为 $1,2,3,4$ 的概率；(2)利用每组的中点值乘以频数，相加后除以100可得结果；(3)根据表格中的数据完善 $2\times2$ 列联表，计算出 $K^2$ 的观测值，再结合临界值表可得结论.

\textbf{详解：}(1) 由频数分布表可知，该市一天的空气质量等级为1的概率为 $\frac{2+16+25}{100}=0.43$，等级为2的概率为 $\frac{5+10+12}{100}=0.27$，等级为3的概率为 $\frac{6+7+8}{100}=0.21$，等级为4的概率为 $\frac{7+2+0}{100}=0.09$.
(2) 一天中到该公园锻炼的人次的平均数为 $\frac{100\times20+300\times35+500\times45}{100}=350$.
(3) $2\times2$ 列联表如下：

|          | 人次$\leq400$ | 人次$>400$ |
|----------|-------------|-----------|
| 空气质量好 | 33          | 37        |
| 空气质量不好 | 22          | 8         |

$K^2=\frac{100\times(33\times8-37\times22)^2}{55\times45\times70\times30}\approx5.820>3.841$，因此，有95%的把握认为一天中到该公园锻炼的人次与该市当天的空气质量有关.
\end{solution}
\end{question}

\begin{question}[points=12]
如图，在长方体 $ABCD-A_1B_1C_1D_1$ 中，点 $E$，$F$ 分别在棱 $DD_1$，$BB_1$ 上，且 $2DE=ED_1$，$BF=2FB_1$．证明：

(1) 当 $AB=BC$ 时，$EF\perp AC$；

(2) 点 $C_1$ 在平面 $AEF$ 内．

\begin{center}
  \includegraphics[width=.38\linewidth]{question-19-1.jpg}
\end{center}
\begin{solution}
\textbf{答案：}证明见解析

\textbf{分析：}(1)根据正方形性质得 $AC\perp BD$，根据长方体性质得 $AC\perp BB_1$，进而可证 $AC\perp$ 平面 $BB_1D_1D$，即得结果；(2)只需证明 $EC_1\parallel AF$ 即可，在 $CC_1$ 上取点 $M$ 使得 $CM=2MC_1$，再通过平行四边形性质进行证明即可.

\textbf{详解：}(1) 因为长方体 $ABCD-A_1B_1C_1D_1$，所以 $BB_1\perp$ 平面 $ABCD$，所以 $AC\perp BB_1$. 因为 $AB=BC$，所以四边形 $ABCD$ 为正方形，所以 $AC\perp BD$. 因为 $BB_1\cap BD=B$，$BB_1,BD\subset$ 平面 $BB_1D_1D$，因此 $AC\perp$ 平面 $BB_1D_1D$. 因为 $EF\subset$ 平面 $BB_1D_1D$，所以 $AC\perp EF$.
(2) 在 $CC_1$ 上取点 $M$ 使得 $CM=2MC_1$，连接 $DM,MF$. 因为 $D_1E=2ED$，$DD_1\parallel CC_1$，$DD_1=CC_1$，所以 $ED=MC_1$，$ED\parallel MC_1$，所以四边形 $DMC_1E$ 为平行四边形，所以 $DM\parallel EC_1$. 因为 $MF\parallel DA$，$MF=DA$，所以四边形 $MFAD$ 为平行四边形，所以 $DM\parallel AF$，所以 $EC_1\parallel AF$，因此 $C_1$ 在平面 $AEF$ 内.
\end{solution}
\end{question}

\begin{question}[points=12]
已知函数 $f(x)=x^3-kx+k^2$．

(1) 讨论 $f(x)$ 的单调性；

(2) 若 $f(x)$ 有三个零点，求 $k$ 的取值范围．
\begin{solution}
\textbf{答案：}(1) 当 $k\leq0$ 时，$f(x)$ 在 $(-\infty,+\infty)$ 上单调递增；当 $k>0$ 时，$f(x)$ 在 $(-\sqrt{\frac{k}{3}},\sqrt{\frac{k}{3}})$ 上单调递减，在 $(-\infty,-\sqrt{\frac{k}{3}})$ 和 $(\sqrt{\frac{k}{3}},+\infty)$ 上单调递增. (2) $(0,\frac{4}{27})$

\textbf{分析：}(1) $f'(x)=3x^2-k$，对 $k$ 分 $k\leq0$ 和 $k>0$ 两种情况讨论即可；(2) $f(x)$ 有三个零点，由(1)知 $k>0$，且 $\begin{cases} f(-\sqrt{\frac{k}{3}})>0 \\ f(\sqrt{\frac{k}{3}})<0 \end{cases}$，解不等式组得到 $k$ 的范围，再利用零点存在性定理加以说明即可.

\textbf{详解：}(1) 由题，$f'(x)=3x^2-k$. 当 $k\leq0$ 时，$f'(x)\geq0$ 恒成立，所以 $f(x)$ 在 $(-\infty,+\infty)$ 上单调递增；当 $k>0$ 时，令 $f'(x)=0$，得 $x=\pm\sqrt{\frac{k}{3}}$，令 $f'(x)<0$，得 $-\sqrt{\frac{k}{3}}<x<\sqrt{\frac{k}{3}}$，令 $f'(x)>0$，得 $x<-\sqrt{\frac{k}{3}}$ 或 $x>\sqrt{\frac{k}{3}}$，所以 $f(x)$ 在 $(-\sqrt{\frac{k}{3}},\sqrt{\frac{k}{3}})$ 上单调递减，在 $(-\infty,-\sqrt{\frac{k}{3}})$ 和 $(\sqrt{\frac{k}{3}},+\infty)$ 上单调递增.
(2) 由(1)知，$f(x)$ 有三个零点，则 $k>0$，且 $\begin{cases} f(-\sqrt{\frac{k}{3}})>0 \\ f(\sqrt{\frac{k}{3}})<0 \end{cases}$，即 $\begin{cases} k^2+\frac{2}{3}k\sqrt{\frac{k}{3}}>0 \\ k^2-\frac{2}{3}k\sqrt{\frac{k}{3}}<0 \end{cases}$，解得 $0<k<\frac{4}{27}$. 当 $0<k<\frac{4}{27}$ 时，$k>\sqrt{\frac{k}{3}}$，且 $f(k)=k^2>0$，所以 $f(x)$ 在 $(\sqrt{\frac{k}{3}},k)$ 上有唯一一个零点；同理 $-k-1<-\sqrt{\frac{k}{3}}$，$f(-k-1)=-k^3-(k+1)^2<0$，所以 $f(x)$ 在 $(-k-1,-\sqrt{\frac{k}{3}})$ 上有唯一一个零点；又 $f(x)$ 在 $(-\sqrt{\frac{k}{3}},\sqrt{\frac{k}{3}})$ 上有唯一一个零点，所以 $f(x)$ 有三个零点. 综上可知 $k$ 的取值范围为 $(0,\frac{4}{27})$.
\end{solution}
\end{question}

\begin{question}[points=12]
已知椭圆 $C:\frac{x^2}{25}+\frac{y^2}{m^2}=1(0<m<5)$ 的离心率为 $\frac{\sqrt{15}}{4}$，$A$，$B$ 分别为 $C$ 的左、右顶点．

(1) 求 $C$ 的方程；

(2) 若点 $P$ 在 $C$ 上，点 $Q$ 在直线 $x=6$ 上，且 $|BP|=|BQ|$，$BP\perp BQ$，求 $\triangle APQ$ 的面积．
\begin{solution}
\textbf{答案：}(1) $\frac{x^2}{25}+\frac{16y^2}{25}=1$ (2) $\frac{5}{2}$

\textbf{分析：}(1) 由椭圆方程可得 $a=5$, $b=m$，根据离心率公式结合已知求得 $m$ 即可；(2) 点 $P$ 在 $C$ 上，点 $Q$ 在直线 $x=6$ 上，且 $|BP|=|BQ|$，$BP\perp BQ$，过点 $P$ 作 $x$ 轴垂线，交点为 $M$，设 $x=6$ 与 $x$ 轴交点为 $N$，可得 $\triangle PMB\cong\triangle BNQ$，可求得 $P$ 点坐标，求出直线 $AQ$ 的方程，根据点到直线距离公式和两点距离公式，即可求得 $\triangle APQ$ 的面积.

\textbf{详解：}(1) 因为 $C:\frac{x^2}{25}+\frac{y^2}{m^2}=1(0<m<5)$，所以 $a=5$, $b=m$. 离心率 $e=\frac{c}{a}=\sqrt{1-\frac{b^2}{a^2}}=\sqrt{1-\frac{m^2}{25}}=\frac{\sqrt{15}}{4}$，解得 $m=\frac{5}{4}$ (负舍). 所以 $C$ 的方程为 $\frac{x^2}{25}+\frac{y^2}{(5/4)^2}=1$，即 $\frac{x^2}{25}+\frac{16y^2}{25}=1$.
(2) 由椭圆方程知 $B(5,0)$，设 $P(x_P,y_P)$，$Q(6,y_Q)$. 过 $P$ 作 $PM\perp x$ 轴于 $M$，过 $Q$ 作 $QN\perp x$ 轴于 $N$，则 $\triangle PMB$ 与 $\triangle BNQ$ 中，$\angle PMB=\angle BNQ=90^\circ$，$BP=BQ$，$\angle PBM+\angle QBN=90^\circ$，$\angle BQN+\angle QBN=90^\circ$，所以 $\angle PBM=\angle BQN$，所以 $\triangle PMB\cong\triangle BNQ$，因此 $PM=BN$, $MB=NQ$. 由 $N(6,0)$ 得 $BN=6-5=1$，所以 $PM=1$，即 $y_P=1$ (由对称性可取正值). 代入椭圆方程得 $\frac{x_P^2}{25}+\frac{16}{25}=1$，解得 $x_P=3$ 或 $x_P=-3$.
① 若 $P(3,1)$，则 $MB=5-3=2$，所以 $NQ=MB=2$，得 $Q(6,2)$. $A(-5,0)$，直线 $AQ$ 的方程为 $\frac{y-0}{2-0}=\frac{x+5}{6+5}$，即 $2x-11y+10=0$. 点 $P$ 到直线 $AQ$ 的距离 $d=\frac{|2\times3-11\times1+10|}{\sqrt{2^2+(-11)^2}}=\frac{5}{\sqrt{125}}=\frac{\sqrt{5}}{5}$，$|AQ|=\sqrt{(6+5)^2+(2-0)^2}=5\sqrt{5}$，所以 $S_{\triangle APQ}=\frac{1}{2}\times5\sqrt{5}\times\frac{\sqrt{5}}{5}=\frac{5}{2}$.
② 若 $P(-3,1)$，则 $MB=5-(-3)=8$，所以 $NQ=8$，得 $Q(6,8)$. $A(-5,0)$，直线 $AQ$ 的方程为 $\frac{y-0}{8-0}=\frac{x+5}{6+5}$，即 $8x-11y+40=0$. 点 $P$ 到直线 $AQ$ 的距离 $d=\frac{|8\times(-3)-11\times1+40|}{\sqrt{8^2+(-11)^2}}=\frac{5}{\sqrt{185}}=\frac{\sqrt{185}}{37}$，$|AQ|=\sqrt{(6+5)^2+(8-0)^2}=\sqrt{185}$，所以 $S_{\triangle APQ}=\frac{1}{2}\times\sqrt{185}\times\frac{\sqrt{185}}{37}=\frac{185}{74}=\frac{5}{2}$.
综上所述，$\triangle APQ$ 的面积为 $\frac{5}{2}$.
\end{solution}
\end{question}

\begin{question}[points=10]
[选修4-4：坐标系与参数方程]

在直角坐标系 $xOy$ 中，曲线 $C$ 的参数方程为 $\begin{cases} x=2-t-t^2 \\ y=2-3t+t^2 \end{cases}$ ($t$ 为参数且 $t\neq1$)，$C$ 与坐标轴交于 $A$，$B$ 两点．

(1) 求 $|AB|$；

(2) 以坐标原点为极点，$x$ 轴正半轴为极轴建立极坐标系，求直线 $AB$ 的极坐标方程．
\begin{solution}
\textbf{答案：}(1) $4\sqrt{10}$ (2) $3\rho\cos\theta-\rho\sin\theta+12=0$

\textbf{分析：}(1) 由参数方程得出 $A,B$ 的坐标，再由两点间距离公式即可得出 $|AB|$ 的值；(2) 由 $A,B$ 的坐标得出直线 $AB$ 的直角坐标方程，再化为极坐标方程即可.

\textbf{详解：}(1) 令 $x=0$，则 $t^2+t-2=0$，解得 $t=-2$ 或 $t=1$ (舍)，则 $y=2+6+4=12$，即 $A(0,12)$. 令 $y=0$，则 $t^2-3t+2=0$，解得 $t=2$ 或 $t=1$ (舍)，则 $x=2-2-4=-4$，即 $B(-4,0)$. 所以 $|AB|=\sqrt{(0+4)^2+(12-0)^2}=4\sqrt{10}$.
(2) 由(1)得 $k_{AB}=\frac{12-0}{0-(-4)}=3$，则直线 $AB$ 的方程为 $y=3(x+4)$，即 $3x-y+12=0$. 由 $x=\rho\cos\theta$, $y=\rho\sin\theta$ 可得，直线 $AB$ 的极坐标方程为 $3\rho\cos\theta-\rho\sin\theta+12=0$.
\end{solution}
\end{question}

\begin{question}[points=10]
[选修4-5：不等式选讲]

设 $a,b,c\in\mathbb R$，$a+b+c=0$，$abc=1$．

(1) 证明：$ab+bc+ca<0$；

(2) 用 $\max\{a,b,c\}$ 表示 $a,b,c$ 中的最大值，证明：$\max\{a,b,c\}\geq\sqrt[3]{4}$．
\begin{solution}
\textbf{答案：}证明见解析

\textbf{分析：}(1) 由 $(a+b+c)^2=a^2+b^2+c^2+2ab+2ac+2bc=0$ 结合不等式的性质，即可得出证明；(2) 不妨设 $\max\{a,b,c\}=a$，由题意得出 $a>0,b<0,c<0$，由 $a^3=a^2\cdot a=\frac{(b+c)^2}{bc}=\frac{b^2+c^2+2bc}{bc}$，结合基本不等式，即可得出证明.

\textbf{详解：}(1) 因为 $(a+b+c)^2=a^2+b^2+c^2+2ab+2ac+2bc=0$，所以 $ab+bc+ca=-\frac{1}{2}(a^2+b^2+c^2)$. 由于 $a,b,c$ 均不为0，则 $a^2+b^2+c^2>0$，故 $ab+bc+ca=-\frac{1}{2}(a^2+b^2+c^2)<0$.
(2) 不妨设 $\max\{a,b,c\}=a$，由 $a+b+c=0$，$abc=1$ 可知 $a>0$, $b<0$, $c<0$. 因为 $a=-b-c$，$a=\frac{1}{bc}$，所以 $a^3=a^2\cdot a=\frac{(b+c)^2}{bc}=\frac{b^2+c^2+2bc}{bc}\geq\frac{2bc+2bc}{bc}=4$，当且仅当 $b=c$ 时取等号，故 $a\geq\sqrt[3]{4}$，即 $\max\{a,b,c\}\geq\sqrt[3]{4}$.
\end{solution}
\end{question}

